#m7
\documentclass[11pt,a4paper]{article}

\usepackage{amsmath}
\usepackage{amssymb}
\usepackage{amsthm}
\usepackage{geometry}
\usepackage{booktabs}

\geometry{left=2.5cm, right=2.5cm, top=2.5cm, bottom=2.5cm}

\newtheorem{theorem}{Theorem}
\newtheorem{lemma}{Lemma}

\begin{document}

\title{Autonomous Multiclass TRM: Spatial Aggregation, Probabilistic Moving Bottlenecks, and Riemann Benchmarks}
\author{Mingchen Yuan}
\date{}
\maketitle

\section{Introduction}
This paper presents a unified 1D aggregated macroscopic traffic flow model that couples foundational fluid-dynamic conservation laws with a probabilistic subclass extension for moving bottlenecks. Following the principle of dimensionality reduction, we aggregate multi-lane dynamics into a cross-sectional continuum. By utilizing a physics-based probabilistic approximation, this model elegantly captures how heavy trucks trap faster passenger cars without relying on computationally expensive, rigid lane-by-lane topological constraints. The framework is strictly designed to satisfy absolute macroscopic benchmarks---including global mass conservation, non-negative densities, and strict jam density bounds---and reliably reproduces standard fluid-dynamic Riemann phenomena such as shockwaves and rarefaction waves, ensuring absolute stability for subsequent numerical simulations and gradient-based optimizations.

\section{Aggregated State Variables and Isomorphic Subclasses}
Let the highway be discretized into $X$ longitudinal cells $x\in\{1, \dots, X\}$ of length $\Delta x$. By aggregating all lanes, the traffic flow is modeled across the entire cross-section. Vehicles travel at discrete speed levels $i \in \{1, \dots, N\}$, corresponding to physical speeds $v_1 < v_2 < \dots < v_N$. 

To phenomenologically capture the moving bottleneck effect without rigid micro-topological bounds, we partition the continuum into three distinct interacting subclasses ($m$):
\begin{itemize}
    \item \textbf{Class A (Trucks):} Heavy vehicles with a free-flow speed limit of $v_A^{ff}$.
    \item \textbf{Class Bf (B-fast):} Freely flowing passenger cars.
    \item \textbf{Class Bs (B-slow):} Passenger cars probabilistically trapped behind trucks.
\end{itemize}

To strictly prevent unphysical spatial overlapping, the absolute physical capacity limit $\rho_{\max}$ is governed by the \textbf{Effective Spatial Occupancy Density}, weighted by the dimensionless Passenger Car Equivalent (PCE) $w^{(m)}$:
\begin{equation}
\Omega_{x} = \sum_{m\in\{A,Bf,Bs\}}\sum_{i=1}^{N} w^{(m)} f_{i,x}^{(m)}
\end{equation}
where $w^{(Bf)} = w^{(Bs)} = w^{(B)}$, ensuring that a passenger car's physical spatial footprint is strictly invariant regardless of its operational state.

\subsection{Dynamic Asymmetric Kinematics for Trapped Cars}
The maximum static speed index accessible to a truck is $i_{thr} = \max\{i : v_i \le v_A^{ff}\}$. Both free cars ($Bf$) and trapped cars ($Bs$) mathematically coexist across the entire velocity spectrum. Because lateral lane-changing is implicitly aggregated in this 1D continuum, we mathematically define the "trapped" state by enforcing an absolute asymmetric acceleration blockade exclusively upon Class Bs. 

To prevent "ghost penetration" (where a trapped car unrealistically accelerates past a severely congested truck), this blockade dynamically tracks the local bottleneck. Let $\kappa_x(t) \le i_{thr}$ represent the \textbf{dynamic bottleneck speed limit} in cell $x$. To avoid numerical diffusion artifacts common in finite volume methods, we define this using an entity threshold $\epsilon_{veh} = 10^{-4}$:
\begin{equation}
\kappa_x(t) = 
\begin{cases} 
\max\{k \in \{1, \dots, i_{thr}\} \mid f_{k,x}^{(A)} > \epsilon_{veh}\}, & \text{if } \sum_{k} f_{k,x}^{(A)} > \epsilon_{veh} \\ 
i_{thr}, & \text{otherwise} 
\end{cases}
\end{equation}

A trapped car is strictly forbidden from accelerating past the current actual physical speed of the truck cluster ahead of it, though normal deceleration towards zero is flawlessly preserved to respect background congestion:
\begin{equation}
\lambda_{i\rightarrow i+1}^{(Bs)} \equiv 0, \quad \forall i \ge \kappa_x(t)
\end{equation}

\section{Baseline Internal Kinematics ($\lambda$)}
For standard kinetic evolutions, velocity transitions strictly occur between adjacent states ($|i-j|\le 1$) with sealed boundary conditions ($\lambda_{1\rightarrow 0} \equiv 0$, $\lambda_{N\rightarrow N+1} \equiv 0$). 

A vehicle's acceleration capability organically decays as it approaches its physical maximum speed, governed by a decay exponent $\delta^{(m)}$ and a base acceleration parameter $a_i^{(m)}$:
\begin{equation}
\lambda_{i\rightarrow i+1}^{(m)} = a_{i}^{(m)}\left(1-\frac{v_{i}}{v_{\max}}\right)^{\delta^{(m)}} \left[1-\exp\left(-\frac{\max(0,\rho_{\max}-\Omega_{x})}{R_{c}}\right)\right]
\end{equation}

To rigorously prevent vehicles from colliding into a jammed downstream, deceleration is governed by a \textbf{Singular Asymptotic Barrier} $\mathcal{B}(\Omega_{x})$. A regularizer $\epsilon \approx 10^{-8}$ completely eradicates division-by-zero anomalies:
\begin{equation}
\mathcal{B}(\Omega_{x}) = \left(\frac{\rho_{\max}}{\max(\epsilon,\rho_{\max}-\Omega_{x})}\right)^{\eta^{(m)}}
\end{equation}
\begin{equation}
\lambda_{i\rightarrow i-1}^{(m)} = \left[ \omega_{0}^{(m)} + \sum_{n}\sum_{k<i}\beta_{i\rightarrow i-1}^{(m,n)}(v_{i}-v_{k}) \frac{w^{(n)} f_{k,x}^{(n)}}{\rho_{\max}} \right] \cdot \mathcal{B}(\Omega_{x})
\end{equation}

\section{Macroscopic Probabilistic Reactions: Capture and Release}
Rather than utilizing rigid topological blocking, we adopt a \textbf{Probabilistic Approximation}. The fundamental logic dictates: the macroscopic probability of a fast vehicle being trapped is strictly proportional to both its kinetic exposure to slower trucks and the overall macroscopic lack of overtaking opportunities.

\subsection{Capture Rate: Probabilistic Entrapment ($Bf \rightarrow Bs$)}
First, we evaluate the normalized kinetic exposure to downstream trucks. The indicator function $\mathbf{1}_{v_i \ge v_k}$ structurally guarantees that cars are only exposed to trucks operating at slower or equal speeds. \textit{(Note: Since $v_i, v_k$ are fixed speed coordinates, this indicator is a constant mask matrix, preserving the continuous differentiability of the state variables during numerical optimization)}:
\begin{equation}
\mathcal{A}_{x}^{(i)} = \sum_{k \le i_{thr}} \mathbf{1}_{v_i \ge v_k} \left[ \omega_0^{(B,A)} + \beta^{(B,A)}(v_i - v_k) \right] \frac{w^{(A)} f_{k,x}^{(A)}}{\rho_{\max}}
\end{equation}

Second, we define the \textbf{Probabilistic Blocking Factor} $P_{block}(x)$. This acts as the macroscopic statistical equivalent of blocked overtaking lanes. The probability of failed overtaking scales exponentially as the total spatial occupancy approaches the physical jam density:
\begin{equation}
P_{block}(x) = \left( \frac{\min[\max(0, \Omega_{x}), \rho_{\max}]}{\rho_{\max}} \right)^{\eta_{block}}
\end{equation}
\textit{Macro-Averaging Logic: If slow trucks account for a significant portion of capacity and overall congestion restricts passing maneuvers (e.g., 30\% of the road is occupied), fast cars face a proportional, mathematically continuous probability of failing to overtake, thereby transitioning into the trapped Bs state.}

The total isomorphic capture rate mathematically multiplies these probabilities, amplified by the singular barrier near capacity limits:
\begin{equation}
\sigma_{x}^{(i)} = \sigma_0^{(B)} \cdot P_{block}(x) \cdot \mathcal{A}_{x}^{(i)} \cdot \mathcal{B}(\Omega_{x})
\end{equation}

\subsection{Release Rate: Dissipation of Bottlenecks ($Bs \rightarrow Bf$)}
Release is the inverse thermodynamic process. Trapped cars organically recover their free-flow status as the localized truck cluster disperses or background congestion drops. Let the dimensionless local truck footprint be $\theta_{x}^{(A)} = \sum_{k \le i_{thr}} \frac{w^{(A)} f_{k,x}^{(A)}}{\rho_{\max}}$.

To prevent high-frequency numerical chattering and ensure parameter consistency for numerical calibration algorithms, the effective bottleneck pressure $\tilde{\mathcal{S}}_x$ perfectly mirrors the capture logic and remains strictly dimensionless ($\tilde{\mathcal{S}}_x \in [0,1]$):
\begin{equation}
\tilde{\mathcal{S}}_x = P_{block}(x) \cdot \theta_{x}^{(A)}
\end{equation}

The release rate follows the exponential decay of this effective pressure, where $R_A$ operates purely as a dimensionless sensitivity scalar:
\begin{equation}
\mu_{x} = \mu_0^{(B)} \exp\left(-\frac{\tilde{\mathcal{S}}_{x}}{R_A}\right)
\end{equation}
\textit{Physical Axiom: If the trucks accelerate away ($\theta_{x}^{(A)} \rightarrow 0$) or the overall congestion clears enabling overtaking ($P_{block} \rightarrow 0$), the macroscopic blocking pressure instantly dissipates ($\tilde{\mathcal{S}}_x \rightarrow 0$). The release rate smoothly surges to its maximum $\mu_0^{(B)}$, organically freeing the trapped vehicles.}

\section{Spatial Advection and Differentiable Godunov Limiter}
All vehicles fundamentally advect forward at their actual instantaneous speed $v_i$. The raw continuous spatial demand flux is modeled as:
\begin{equation}
\Psi_{i,x\rightarrow x+1}^{(m)} = v_{i} f_{i,x}^{(m)} \cdot \left[1 - \exp\left(-\frac{\max(0, \rho_{\max} - \Omega_{x+1})}{R_{supply}}\right)\right]
\end{equation}

Evaluating continuous exponential filters within an explicit discrete time step $\Delta t$ inherently risks mass overshooting. To mathematically guarantee absolute physical boundedness without breaking computational graph continuity (crucial for future gradient-based numerical optimizations like Automatic Differentiation), we introduce a fully differentiable continuous Godunov flux limiter $\alpha$. Let the total spatial demand be $\mathcal{D}_{x\rightarrow x+1} = \frac{\Delta t}{\Delta x} \sum_{n, j} w^{(n)} \Psi_{j,x\rightarrow x+1}^{(n)}$. We integrate a regularized machine tolerance $\epsilon_{tol} \approx 10^{-12}$ to safely eliminate division-by-zero singularities without relying on non-differentiable \texttt{if-else} branches:
\begin{equation}
\alpha_{x\rightarrow x+1} = \min \left[ 1, \frac{\max(0, \rho_{\max} - \Omega_{x+1})}{\mathcal{D}_{x\rightarrow x+1} + \epsilon_{tol}} \right]
\end{equation}
The strictly bounded, non-negative Eulerian transfer flux is uniformly scaled: $\Phi_{i,x\rightarrow x+1}^{(m)} = \alpha_{x\rightarrow x+1} \cdot \Psi_{i,x\rightarrow x+1}^{(m)}$.

\section{Unified 1D Governing PDEs}
Because lanes are aggregated, the governing Partial Differential Equations strictly reduce to 1D space plus the discrete velocity spectrum. Capture and release reactions occur isomorphically at the exact speed index $i$.

\textbf{Class A (Trucks):}
\begin{equation}
\frac{\partial f_{i,x}^{(A)}}{\partial t} + \frac{\partial \Phi_{i,x}^{(A)}}{\partial x} = \sum_{j\in\{i-1,i+1\}} \left( \lambda_{j\rightarrow i}^{(A)} f_{j,x}^{(A)} - \lambda_{i\rightarrow j}^{(A)} f_{i,x}^{(A)} \right)
\end{equation}

\textbf{Class Bf (Free-flowing Cars):}
\begin{equation}
\begin{aligned}
\frac{\partial f_{i,x}^{(Bf)}}{\partial t} + \frac{\partial \Phi_{i,x}^{(Bf)}}{\partial x} &= 
\sum_{j\in\{i-1,i+1\}} \left( \lambda_{j\rightarrow i}^{(Bf)} f_{j,x}^{(Bf)} - \lambda_{i\rightarrow j}^{(Bf)} f_{i,x}^{(Bf)} \right) \\
&\quad + \mu_{x} f_{i,x}^{(Bs)} - \sigma_{x}^{(i)} f_{i,x}^{(Bf)}
\end{aligned}
\end{equation}

\textbf{Class Bs (Trapped Cars):}
\begin{equation}
\begin{aligned}
\frac{\partial f_{i,x}^{(Bs)}}{\partial t} + \frac{\partial \Phi_{i,x}^{(Bs)}}{\partial x} &= 
\sum_{j\in\{i-1,i+1\}} \left( \lambda_{j\rightarrow i}^{(Bs)} f_{j,x}^{(Bs)} - \lambda_{i\rightarrow j}^{(Bs)} f_{i,x}^{(Bs)} \right) \\
&\quad + \sigma_{x}^{(i)} f_{i,x}^{(Bf)} - \mu_{x} f_{i,x}^{(Bs)}
\end{aligned}
\end{equation}

\section{Mathematical Proof of Absolute Macroscopic Benchmarks}
\begin{theorem}[Strict Conservation and Positivity Bounds]
The macroscopic PDE system and its discrete numerical mappings mathematically guarantee that: (1) Total global mass is unconditionally conserved; (2) All density states remain strictly non-negative given the Courant-Friedrichs-Lewy (CFL) constraint; (3) The physical jam density $\rho_{\max}$ is never exceeded.
\end{theorem}
\begin{proof}
\textbf{1. Mass Conservation:} Summing the PDEs over the discrete velocity spectrum yields absolute zero-sum cancellations. The internal isomorphic reactions sum to zero: $\sum_i (\mu_x f_{i,x}^{(Bs)} - \sigma_x^{(i)} f_{i,x}^{(Bf)}) + \sum_i (\sigma_x^{(i)} f_{i,x}^{(Bf)} - \mu_x f_{i,x}^{(Bs)}) \equiv 0$. The tridiagonal internal speed changes perfectly cancel out. The temporal derivative reduces exclusively to the boundary spatial flux divergence, strictly satisfying the continuum mass conservation law.

\textbf{2. Positivity Preservation:} The analytical exponential integration used for reactions (Phase 1) structurally bounds components strictly above zero. In Phase 3, provided that the explicit integration strictly respects the \textbf{CFL condition} ($v_{\max} \frac{\Delta t}{\Delta x} \le 1$), the outward advective mass fraction never exceeds the available inventory in the source cell, mathematically guaranteeing $f_{i,x}^{(m)}(t+\Delta t) \ge 0$.

\textbf{3. Jam Density Limit:} During the spatial advection update, the differentiable Godunov limiter $\alpha$ strictly truncates the total advection flux such that the integrated mass influx mathematically cannot exceed the remaining downstream void space $\max(0, \rho_{\max} - \Omega_{x+1})$. Therefore, $\Omega_x \le \rho_{\max}$ is permanently preserved.
\end{proof}

\section{Standard Traffic Scenarios: Riemann Problems Validation}
To validate the model's hyperbolic fluid-dynamical integrity in subsequent numerical experiments, the formulation accurately resolves classical Riemann Problems without non-physical oscillations:
\begin{itemize}
    \item \textbf{Shockwave Generation (Traffic Jams):} Initializing upstream with free-flow states ($\Omega \ll \rho_{\max}$) and downstream with heavy congestion ($\Omega \to \rho_{\max}$). The Godunov limiter mathematically forces $\Phi \to 0$ at the interface. Simultaneously, $\mathcal{B}(\Omega_{x}) \to \infty$, triggering instant momentum thermalization ($\lambda_{i\to i-1} \to \infty$). This seamlessly creates a stable, backward-propagating shockwave that strictly satisfies macroscopic Rankine-Hugoniot jump conditions.
    \item \textbf{Rarefaction Wave (Queue Discharge):} Initializing upstream with jammed states and downstream with free-flow states. As the downstream bottleneck clears, $\mathcal{B}(\Omega_x) \rightarrow 1$. The continuous acceleration operator ($\lambda_{i \to i+1}$) smoothly diffuses the leading edge of the platoon, correctly simulating a continuous fan-like expansion wave indicative of physical acceleration.
\end{itemize}

\section{Operator Splitting: Strict Causal Resolution}
To evaluate the multiscale dynamics efficiently and prevent "ghost penetration" (where trapped cars advect past trucks before their speeds resolve), the system decouples into a strict 3-phase causal sequence per $\Delta t$:

\begin{table}[h]
\centering
\begin{tabular}{llcl}
\toprule
\textbf{Phase} & \textbf{Operator} & \textbf{Stiffness} & \textbf{Numerical Solver} \\
\midrule
1 & Capture \& Release ($\sigma, \mu$) & \textbf{Extremely High} & Exact Analytical Matrix Exponential \\
2 & Internal Speed Changes ($\lambda$) & \textbf{Extremely High} & Dynamic Algebraic Projection + Thomas \\
3 & Spatial Advection ($\Phi$) & Low & Explicit Euler + Differentiable Limiter \\
\bottomrule
\end{tabular}
\end{table}

Let intermediate states be denoted by $(*, **)$. 

\textbf{Phase 1: Analytical Capture \& Release ($t \rightarrow t^*$)} \\
With advection frozen, $\Omega_x$ is invariant. The rate sum $S^{(i)} = \sigma_x^{(i)} + \mu_x$ is a constant scalar. To prevent IEEE 754 division-by-zero singularities when $S^{(i)} \to 0$, we utilize the entire mathematical function $\phi(z) = (1 - e^{-z})/z$ (evaluated securely via \texttt{-expm1(-z)/z}). Let $F_{i,x} = f_{i,x}^{(Bf)}(t) + f_{i,x}^{(Bs)}(t)$; the exact analytical update is evaluated efficiently without numerical truncation errors:
\begin{equation}
f_{i,x}^{(Bf), *} = f_{i,x}^{(Bf)}(t) \exp(-S^{(i)} \Delta t) + \mu_{x} F_{i,x} \Delta t \cdot \phi(S^{(i)} \Delta t)
\end{equation}
\begin{equation}
f_{i,x}^{(Bs), *} = F_{i,x} - f_{i,x}^{(Bf), *}
\end{equation}

\textbf{Phase 2: Instant Slowdown and Kinematics ($t^* \rightarrow t^{**}$)} \\
Before allowing spatial advection, newly trapped fast cars must instantaneously drop to the dynamic bottleneck speed to avoid spatial penetration. Using the dynamic limit $\kappa_x(t)$ defined in Equation (2), an Exact Algebraic Projection collapses all over-speed trapped density:
\begin{equation}
f_{\kappa_x(t),x}^{(Bs), *} \leftarrow f_{\kappa_x(t),x}^{(Bs), *} + \sum_{j > \kappa_x(t)} f_{j,x}^{(Bs), *} ; \quad f_{j,x}^{(Bs), *} \leftarrow 0 \quad (\forall j > \kappa_x(t))
\end{equation}
Subsequently, the unconditionally stable implicit Thomas algorithm efficiently resolves the remaining regular kinematic transitions ($\lambda$). 

\textbf{Phase 3: Spatial Advection ($t^{**} \rightarrow t+\Delta t$)} \\
Because Phases 1 and 2 execute \textit{prior} to Phase 3, trapped vehicles are mathematically anchored to the moving bottleneck's actual velocity. Finally, applying the explicit upstream-formulated Godunov limiter, conditioned by the CFL requirement ($\frac{v_{\max}\Delta t}{\Delta x} \le 1$), guarantees that spatial density propagation rigorously obeys macroscopic physical limits:
\begin{equation}
f_{i,x}^{(m)}(t+\Delta t) = f_{i,x}^{(m), **} + \frac{\Delta t}{\Delta x} \left( \Phi_{i,x-1\rightarrow x}^{(m), **} - \Phi_{i,x\rightarrow x+1}^{(m), **} \right)
\end{equation}

\end{document}